\section{A Continuation That Should Not Be Discarded}
\label{sec:motivation}
Consider a robot that repeatedly visits three checkpoints and a charger on a
$5\times5$ grid. Charging requires a dwell period, and a request monitor
records whether a charging visit is owed. These duties come from the public
Patrolling specification, corpus~277~\cite{kindartifact}. We use its unchanged
source as the new model and construct an old model that additionally fixes
the displayed layout and keeps all checkpoints open. Removing those old ENV
restrictions leaves system requirements, types and monitors unchanged.
Checkpoint positions still remain fixed from one round to the next, but a
distant checkpoint may now close. This update scenario and reuse policy
were constructed after diagnosis.

\subsection{The retained entry and two safety rules}
An allowed old-model step reaches entry $r$: the robot is at the origin
$O=(0,0)$, all checkpoints are open, charging is unavailable, the dwell
counter $c$ is zero and the charging request is pending. Its retained route phase
is ready to visit checkpoint $A$. The checkpoints are at
$(0,4),(2,4),(4,4)$ and the charger at $H=(4,0)$
(Fig.~\ref{fig:patrol-grid}).
% Legal source-policy route only. No off-assumption completion is depicted.
% Uses the main document's TikZ, arrows.meta, sourceblue/fliporange colors.
\begin{figure}[!ht]
\centering
\begin{tikzpicture}[
 x=10mm,y=10mm,
 >={Latex[length=1.5mm,width=1.05mm]},
 every node/.style={font=\small,inner sep=1pt},
 patrol/.style={sourceblue,line width=.8pt,<->,shorten <=1.7mm,shorten >=1.7mm},
 charge/.style={fliporange,line width=.9pt,->,shorten <=1.7mm,shorten >=1.7mm},
 checkpoint/.style={draw=sourceblue,fill=softblue,rectangle,
   minimum width=4.7mm,minimum height=4.2mm,font=\small\bfseries},
 standoff/.style={draw=sourceblue,fill=white,circle,
   line width=.8pt,minimum size=3.4mm,inner sep=0pt}
]
\path[use as bounding box] (-.65,-.65) rectangle (7.85,4.45);

% The grid points, not diagonals or geometric shortcuts, are the source states.
\draw[step=1,black!17,line width=.3pt] (0,0) grid (4,4);
\foreach \x in {0,...,4} {
  \foreach \y in {0,...,4} {
    \fill[black!40] (\x,\y) circle[radius=.35pt];
  }
  \node[anchor=north,text=black!75] at (\x,-.27) {$\x$};
}
\foreach \y in {0,...,4} {
  \node[anchor=east,text=black!75] at (-.28,\y) {$\y$};
}
\node[anchor=west] at (4.30,-.35) {$x$};
\node[anchor=south] at (-.36,4.15) {$y$};

% O to the first standoff, and the return to O after the checkpoint visits.
\foreach \y in {0,1} {
  \draw[patrol] (0,\y)--(0,{\y+1});
}
% Row 2 connects the standoffs; these edges are reused on return journeys.
\foreach \x in {0,1,2,3} {
  \draw[patrol] (\x,2)--({\x+1},2);
}
% Each visit approaches through row 3 and retreats along the same unit edges.
\foreach \x in {0,2,4} {
  \foreach \y in {2,3} {
    \draw[patrol] (\x,\y)--(\x,{\y+1});
  }
}
% Charging: O along row 0 to H; after dwelling, H to (4,1) to (4,2).
\foreach \x in {0,1,2,3} {
  \draw[charge] (\x,0)--({\x+1},0);
}
\foreach \y in {0,1} {
  \draw[charge] (4,\y)--(4,{\y+1});
}

\foreach \x in {0,2,4} {\node[standoff] at (\x,2) {};}
\node[checkpoint] at (0,4) {$A$};
\node[checkpoint] at (2,4) {$B$};
\node[checkpoint] at (4,4) {$C$};
\node[draw=sourceblue,fill=white,circle,minimum size=4.7mm] at (0,0) {$O$};
\node[draw=fliporange,fill=softorange,rectangle,minimum size=4.7mm] at (4,0) {$H$};

\node[anchor=north west,align=left,text width=29mm] at (4.80,4.30)
 {\textcolor{sourceblue}{\textbf{Visit cycle}}\\
 $A$, $B$, $C$ in order;\\
 return to row 2\\
 after each visit.\\[3pt]
 \textcolor{fliporange}{\textbf{Charging leg}}\\
 After $C$, via row 2\\
 and $O$ to $H$;\\
 dwell until $c=3$.\\
 Return to $(4,2)$,\\
 then $A$'s standoff.};
\node[standoff] at (4.98,.15) {};
\node[anchor=west] at (5.27,.15) {standoff};
\end{tikzpicture}
\caption{Patrol's legal source route uses unit cardinal moves.
Before each checkpoint visit, wait at its circled standoff while the next
blocked flag $d_i'$ is true; otherwise approach, visit and retreat.
The charging leg restarts the visit cycle.}
\label{fig:patrol-grid}
\end{figure}


Write $p$ for the robot position and $z$ for charging availability; primes denote next values. A checkpoint $i$ has
blocked flag $d_i$ and position $t_i$. Two SYS safety duties will conflict:
\begin{equation}
\begin{aligned}
\text{Dwell:}\quad &(p=H\land c<3)\Rightarrow p'=H,\\
\text{Avoidance:}\quad &d_i\Rightarrow p'\ne t_i'.
\end{aligned}
\label{eq:patrol-conflict}
\end{equation}
Avoidance combines the \emph{current} blocked flag with the \emph{next}
checkpoint position. Source-legal inputs keep positions fixed and cannot
change a flag while the robot is within one cell of that checkpoint.
The complete ENV justice list requires only that each of the three
checkpoints reopen infinitely often. There is \emph{no charging-availability
justice}: availability is governed by the safety rule $c=3\Rightarrow z'$,
using the current counter.

\subsection{Responding now is not enough}
At $r$, staying at $O$ is safe even for an arbitrary typed next input:
the robot is not dwelling and every current blocked flag is false.
It also has a source policy that visits the checkpoints, waits for reopening
when needed, and charges repeatedly; Section~\ref{sec:patrol} proves this
using the displayed route.

The helper encoding adds an internal ENV Boolean $b$, initially false,
that can change to true once. This is a game input, not the physical update.
Table~\ref{tab:timing} shows the challenge: the first change releases ENV
safety immediately while SYS still owes one safe response. The strict
meaning of that duty is defined in Section~\ref{sec:formal}.
\begin{table}[!ht]
\caption{One edge separates the releases for the complete literal relations
in Eq.~\eqref{hpb:literal}. ``Any typed'' includes source-forbidden inputs.
The last column uses the existing corrected guard.}
\label{tab:timing}\centering\small
\begin{tabular}{@{}llll@{}}\toprule
Helper & Legal ENV input & SYS duty & Corrected duty\\\midrule
$0\to0$ & Source-legal & Source & Source\\
$0\to1$ & Any typed & Source & None\\
$1\to1$ & Any typed & None & None\\\bottomrule
\end{tabular}
\end{table}


Table~\ref{tab:patrol-entry} shows how a legal source prefix reaches the
point where that response fails.
ENV keeps the checkpoints open and charging unavailable until the first
arrival at $H$, then legally closes distant checkpoint $A$ on the next round.
\begin{table}[!ht]
\caption{Milestones of Patrol's legal prefix. Reaching $H$ may take
several rounds; the next row is one dwell edge later. At $x$ the request
has cleared but the dwell is unfinished.}
\label{tab:patrol-entry}\centering\small
\begin{tabular}{@{}L{32mm}ccL{23mm}@{}}\toprule
State / event & $p$ & $c$ & Request\\\midrule
Retained entry $r$ & $O$ & 0 & pending\\
First arrival at $H$ & $H$ & 0 & pending\\
Next legal dwell; $A$ closes ($x$) & $H$ & 1 & cleared (idle)\\\bottomrule
\end{tabular}
\end{table}

Clearing the request records the visit; it does not finish the dwell duty,
which remains active until the counter reaches three. From $x$, a first-change input places the already blocked checkpoint at
$H$: Dwell requires the robot to stay there and Avoidance forbids it.
That input belongs to the translated game; moving a checkpoint is not
source-legal. A different safe strategy cannot avoid the issue: never
reaching $H$ leaves the request pending forever, whereas first reaching
$H$ permits the same legal next round and conflict
(Fig.~\ref{fig:patrol-motivation}). The proof will cover every safe strategy.
The question is whether a winning policy can \emph{keep} a response available,
not only whether the retained state has one now.
\begin{figure}[!ht]
\centering
\begin{tikzpicture}[>=Latex,x=1mm,y=1mm,
every node/.style={font=\small,align=center,inner sep=1.5mm},
wide/.style={draw=sourceblue,fill=softblue,rounded corners=2pt,text width=77mm},
branch/.style={draw,rounded corners=2pt,text width=34mm},
phase/.style={fill=white,inner sep=.4mm}]
\node[wide] (start) at (42,72)
 {Post-update $r$: fixed layout,\\all checkpoints open, no charging.\\
 Can answer every typed first-change input safely.\\
 Consider any safe SYS strategy.};
\node[branch,draw=fliporange,fill=softorange] (never) at (20,51)
 {\textbf{Never reach $H$}\\Counter zero; request pending.\\ENV fair; SYS justice fails.};
\node[branch,draw=sourceblue,fill=softblue] (reach) at (64,51)
 {\textbf{First reach $H$}\\Counter zero.\\Request still pending.};
\draw[->,thick,sourceblue] (start.south west)++(8,0)--(never.north);
\draw[->,thick,sourceblue] (start.south east)++(-8,0)--(reach.north);
\node[wide] (closed) at (42,29)
 {State $x$: distant $A$ closed; robot stays at $H$.\\
 Counter one: \textbf{request cleared, dwell unfinished}.\\
 Some typed first-change input has no safe response.};
\draw[->,thick,sourceblue] (reach.south)--(closed.north);
\node[wide,draw=fliporange,fill=softorange] (flip) at (42,7)
 {Encoding-only first change moves blocked $A$ to $H$.\\
 Dwell requires stay; Avoidance requires leave.\\\textbf{No safe response.}};
\draw[->,thick,fliporange,dashed] (closed.south)
 --node[phase,right] {$b:0\to1$} (flip.north);
\end{tikzpicture}
\caption{All-response loss from Patrol's source-winning entry. Solid arrows use source-legal inputs
while $b=0$. The dashed first-change edge follows the legal edge that
clears the request. The two branches cover every safe strategy.}
\label{fig:patrol-motivation}
\end{figure}

